% !TeX spellcheck = es_ES

\section*{Abstract}

The arrival of next-generation massive astrophysical surveys, such as the \textit{Euclid} space telescope and the \textit{Dark Energy Spectroscopic Instrument} (\textit{DESI}), presents the unprecedented challenge of exploiting large volumes of highly heterogeneous data. This work introduces, optimizes, and validates \textbf{AstroPT MM}, a bimodal autoregressive foundation model of $\sim 100.75$ million parameters pretrained in a self-supervised manner. The model unifies the resolved spatial morphology of optical/infrared images from \textit{Euclid} (DR1) with the fine spectroscopic kinematics of spectra from \textit{DESI} (DESI DR1) using a continuous round-robin block token mixing scheme. The curated training catalog consists of $137\,661$ high-quality image-spectrum pairs.

The evaluation of the joint latent manifold using linear probing and multi-layer perceptrons reveals that the unified representation space systematically outperforms isolated unimodal models and task-specific supervised baselines (ResNet18 and Conv1D+Transformer). \textit{AstroPT MM} achieves good determination coefficients in physical parameter estimation, including stellar mass ($\log M_*$, $R^2 = 0.9339$) and star formation rate ($\log\text{SFR}$, $R^2 = 0.7487$). Furthermore, continuous autoregressive pretraining proves highly competitive against traditional contrastive alignment schemes (\textit{AstroCLIP}), demonstrating that hybrid contrastive-autoregressive architectures degrade the physical representation of individual flux patches due to intermediate bottlenecks.

Finally, we document the utility of the bimodal space in unsupervised downstream applications, including k-NN similarity search for rare objects (dual active galactic nuclei), and generative latent diffusion models for image-conditioned spectral synthesis. This work establishes the methodological foundations for large-scale value-added catalogs and the computational scaling of foundation models via advanced architectures.


\section*{Resumen}

La llegada de cartografiados astrofísicos masivos de nueva generación, como el telescopio espacial \textit{Euclid} y el \textit{Dark Energy Spectroscopic Instrument} (\textit{DESI}), presenta el desafío sin precedentes de explotar grandes volúmenes de datos altamente heterogéneos. En este trabajo se presenta, optimiza y valida \textbf{AstroPT MM}, un modelo fundacional bimodal y autorregresivo de $\sim 100.75$ millones de parámetros preentrenado de forma autosupervisada. El modelo unifica la morfología espacial resuelta de imágenes ópticas/infrarrojas de \textit{Euclid} (DR1) con la cinemática espectroscópica fina de espectros de \textit{DESI} (DESI DR1) a través de un esquema de mezcla de tokens por bloques (\textit{token mixing}) continuo. El catálogo curado de entrenamiento consta de $137\,661$ pares imagen-espectro de alta calidad científica.

La evaluación del colector latente conjunto (\textit{Joint}) mediante regresores lineales y perceptrones multicapa (\textit{probing}) revela que el espacio unificado supera de forma sistemática a los modelos unimodales aislados y a los baselines supervisados específicos de referencia (ResNet18 y Conv1D+Transformer). AstroPT MM alcanza buenos coeficientes de determinación en la estimación de la masa estelar ($\log M_*$, $R^2 = 0.9339$) y en la tasa de formación estelar ($\log\text{SFR}$, $R^2 = 0.7487$). Asimismo, el preentrenamiento autorregresivo continuo resulta altamente competitivo frente a los esquemas contrastivos tradicionales de alineamiento global (\textit{AstroCLIP}), demostrando que las arquitecturas híbridas contrastivas-autorregresivas degradan la representatividad física de los parches individuales de flujo al imponer cuellos de botella rígidos a mitad de la red. 

Por último, se documenta la utilidad del espacio bimodal en tareas no supervisadas, incluyendo motores de búsqueda k-NN para objetos raros (núcleos activos de galaxias duales) y modelos generativos de difusión latente para síntesis espectral continua condicionada a imágenes. Este trabajo establece las bases metodológicas para la creación de catálogos de valor añadido de gran escala y el escalamiento computacional de modelos fundacionales mediante arquitecturas avanzadas.